\scanpage{46}
7. \deleted{\unsure{Symétrie|Symétries}} Involution canonique.

\uline{Prop} Soit $B$ Algèbre quadratique sur $\mathcal O_S$.
\margin{\rule{0.35pt}{3\baselineskip}}Alors l'hom $\mathcal O_S$-linéaire
\[
\sigma:u\longmapsto\deleted{\unsure{\mathrm{tr}\,u\,1_B}}=(\mathrm{tr}\,u)e_B-u \tag*{(84)}
\]
est un \uline{automorphisme} d'algèbres involutif,
caractérisé par la condition
\[
u+\sigma u=\mathrm{tr}\,u\,e_B \tag*{(85)}
\]
% margin note written sideways
\margin{On \unsure{appellera|appelle|veut} quasi-réflexion d'un module $B$, rel. à un sous-module $B'$ tel que $B/B'\simeq\check L$ inversible, tt endomorphisme \deleted{\unsure{par}} de la forme $\mathrm{id}+v$, où $v\in L\otimes V\subset\check V\otimes V$ est tel que $\mathrm{tr}\,v=-2$.}

On notera que la donnée (pour $B$ un $\mathcal O_S$-Module
fixé avec $e_B$ \inserted{section} \deleted{\unsure{partout}} non nulle en chaque fibre)
la donnée de l'involution $\sigma$ équivaut à celle
de la forme \inserted{$\underline b=$} $\mathrm{tr}$\deleted{\unsure{$b$}} (section de $2F$), \deleted{\unsure{et nous}}
Moyennant cette correspondance, les $\underline b$
correspondent \deleted{aux} exact\textsuperscript{t} aux \inserted{quasi-}\uline{réflexions} $\sigma$
du $\mathcal O_S$-Module $B$ telles que $\sigma e_B=e_B$, i.e.
admettant $\mathcal O_S\cdot e_B$ comme droite fixe\unsure{/} \deleted{\unsure{soit}}

% slanted
\margin{\deleted{\unsure{Sauf erreur|Sauf en faisant|déf en faisant} si \unsure{on est partout|au point} \unsure{de car. résiduelle|car si résiduellement} \unsure{char|de} 2, $\mathrm{tr}\,\unsure{v|u}$ n'est pas \unsure{forcément nulle|partout nulle} sur les fibres résiduelles\ldots}}
\deleted{Cette réflexion \unsure{n'est effet|n'est|n'est effect.}, un endomorphisme
ayant $\mathcal O_S e_B$ comme droite de pts fixes se met
sous la forme générale}
\[
\deleted{\sigma u=u+a\deleted{\unsure{\otimes}}\pi_1(u)\qquad\text{i.e. }\sigma=\mathrm{id}_B+\deleted{\unsure{a\otimes}}\pi_1\otimes a} \tag*{\deleted{\unsure{( }}}
\]
\deleted{où $a\in\Gamma(B\otimes L\simeq V=\check B)$ et \deleted{\unsure{où}} (et $\pi_1\in\Gamma(\check B\otimes\check L\simeq\underline{\mathrm{Hom}}(B,\check L)\subset B)$ \inserted{pour mém\textsuperscript{re}}}
\[
\begin{aligned}
\deleted{\sigma^2}&\deleted{=\mathrm{id}_V+2\pi_1\otimes a+\deleted{\pi_1(a)}\inserted{\langle a,\pi_1\rangle}\pi_1\otimes a}\\
&\deleted{=\mathrm{id}_V+\pi_1\otimes a(2+\langle a,\pi_1\rangle)}
\end{aligned}
\]
\deleted{donc $\sigma$ est involutive \deleted{\unsure{ssi}} et \deleted{\unsure{part|partout}} $\ne\mathrm{id}$ en tte
fibre i.e. est une réflexion propre, ssi $a$ \deleted{est} $\ne0$
en tte fibre, \uline{et} $\langle a,\pi_1\rangle\deleted{\unsure{a}}=-2$. \deleted{\unsure{voir}} Identifions}
% archive sheet number: 45
